\section{Zadanie 11}
Niech \((X, Y)\) będzie punktem wylosowanym jednostajnie ze zbioru
\(\{(x, y) \mid x^2+y^2 \leqslant 1, x \geqslant 0\}\)
(czyli jednostkowego półkola).
Jaki rozkład ma zmienna \(Y/X\)?
\subsection*{Rozwiązanie}
Poprawmy treść zadania i zamieńmy warunek \(x\geqslant0\) na \(x>0\),
aby uniknąć dzielenia przez zero (\(x=0\) i tak jest zdarzeniem
miary zero).

Zauważmy, że wylosowanie punktu jednostajnie z półkola jest równoważne
z wylosowaniem jednostajnie kąta i (już niejednostajnie) odległości
od \((0, 0)\).

Nazwijmy nasze półkole 
\(S = \{(x, y) \mid x^2+y^2 \leqslant 1, x > 0\}\)

\[
    F(z) = \int_{S} [Y/X \leqslant z]~ds
    = \int_{-\frac{\pi}2}^{\frac{\pi}2}
    \int_{0}^{1} 2r[Y/X \leqslant z]~dr~d\alpha
\]

Zdarzenie \([Y/X \leqslant z]\) jest niezależne od \(r\)

\[
    F(z)
    = \int_{-\frac{\pi}2}^{\frac{\pi}2}[Y/X \leqslant z] \int_{0}^{1}
    \frac{2r}{\pi}~dr~d\alpha
    = \int_{-\frac{\pi}2}^{\frac{\pi}2}[Y/X \leqslant z]
    \frac{1}{\pi}~d\alpha
    = \int_{-\frac{\pi}2}^{\frac{\pi}2}[\tan \alpha \leqslant z]
    \frac{1}{\pi} ~d\alpha
\]
\[
    = \int_{-\frac{\pi}2}^{\tan^{-1}z} \frac{d\alpha}{\pi}
    = \frac{1}{\pi} \left(\tan^{-1} (z) + \frac{\pi}2\right)
\]
Jest to dystrybuanta rozkładu Cauchy'ego

\subsection*{Odpowiedź}
\[
\boxed{\frac{Y}{X} \sim Cauchy (0,1)}
\]
